%------------------------------ Ejercicio 3 -------------------------------
\setcounter{eje}{2}
\begin{ejer}
Suponga que $U$ y $V$ son conformemente equivalentes. Pruebe que si $U$ es simplemente conexo, entonces también lo es $V$. Note que esta conclusión sigue siendo válida si simplemente asumimos que existe una biyección continua entre $U$ y $V$.
\end{ejer}
\begin{sol}
Sea $f:U\to V$ una equivalencia conforme. Entonces $f$ es una
biyección continua y, como $f^{-1}$ también es holomorfa, se sigue que
$f^{-1}
\colon V\to U$ es continua. Por tanto, $f$ es un homeomorfismo.

Supongamos que $U$ es simplemente conexo y sean $\gamma_0,\gamma_1 \colon [a,b] \to V$ dos curvas en $V$ con $\gamma_0(a) = \gamma_1(a)$ y $\gamma_0(b)=\gamma_1(b)$.
Definimos las curvas
\[
\sigma_0\coloneq f^{-1}\circ\gamma_0,
\qquad
\sigma_1\coloneq f^{-1}\circ\gamma_1
\]
Entonces $\sigma_0,\sigma_1 \colon [a,b] \to U$ son curvas con
\[
\sigma_0(a)
=
f^{-1}(\gamma_0(a))
=
f^{-1}(\gamma_1(a))
= \sigma_1(a)
\qquad\text{y}\qquad
\sigma_0(b)
=
f^{-1}(\gamma_0(b))
=
f^{-1}(\gamma_1(b))
= \sigma_1(b)
\]
Como $U$ es simplemente conexo, existe una familia de curvas
$\set*{\sigma_s\colon[a,b]\to U}_{0\le s\le1}$ con $\sigma_s(t)$ conjuntamente continua en $s$ y $t$, de tal modo que $\sigma_s(t)\big|_{s=0}=\sigma_0(t)$ y $\sigma_s(t)\big|_{s=1}=\sigma_1(t)$ para todo $t\en[a,b]$ y con extremos $\sigma_s(a) = \sigma_0(a)$ y $\sigma_s(b) = \sigma_0(b)$ para todo $0\le s \le 1$.

Ahora definimos, para cada $0\le s\le1$, la curva $\gamma_s \coloneq f\circ\sigma_s$ en $V$.
Es elemental probar que $\set{\gamma_s}_{0\le s\le1}$ define una homotopía entre
$\gamma_0$ y $\gamma_1$:
En efecto, para todo $t\en[0,1]$ tenemos:
\[
\gamma_s(t)\big|_{s=0}
=
f(\sigma_0(t))
=
f(f^{-1}(\gamma_0(t)))
=
\gamma_0(t)
\quad\text{y}\quad
\gamma_s(t)\big|_{s=1}
=
f(\sigma_1(t))
=
f(f^{-1}(\gamma_1(t)))
=
\gamma_1(t)
\]
asimismo, para todo $0\le s \le 1$:
\[
\gamma_s(a)=f(\sigma_s(a)) = f(\sigma_0(a))=\gamma_0(a)
\qquad\text{y}\qquad
\gamma_s(b)=f(\sigma_s(b)) = f(\sigma_0(b))=\gamma_0(b)
\]

Finalmente, como $f$ es continua y $\sigma_s(t)$ es conjuntamente
continua en $(s,t)$, se sigue que $\gamma_s(t)$ también es
conjuntamente continua en $(s,t)$.
Por tanto, $\gamma_0$ y $\gamma_1$ son homotópicas en $V$.
Como las curvas elegidas fueron arbitrarias, concluimos que $V$
es simplemente conexo.

Notemos que en esta demostración solo usamos el hecho de $f$ es homeomorfismo. Más aún, el teorema de invarianza del dominio garantiza que toda biyección continua $f \colon U \to V$ será abierta, y por tanto, un homeomorfismo. Luego, para esta propiedad basta con suponer que $f$ es biyección continua.
\end{sol}
